\documentclass[10pt,a4paper]{amsart}
\usepackage{amsmath,amssymb,amsthm,graphicx}
\usepackage{verbatim}
\newenvironment{reference}{\comment}{\endcomment}
%\newenvironment{reference}{}{}
\newenvironment{slogan}{\comment}{\endcomment}
\newenvironment{history}{\comment}{\endcomment}

% For commutative diagrams we use Xy-pic
\usepackage[all]{xy}

% We use 2cell for 2-commutative diagrams.
\xyoption{2cell}
% Theorem environments.
%
\theoremstyle{plain}
\newtheorem{theorem}[subsection]{Theorem}
\newtheorem{proposition}[subsection]{Proposition}
\newtheorem{lemma}[subsection]{Lemma}

\theoremstyle{definition}
\newtheorem{definition}[subsection]{Definition}
\newtheorem{example}[subsection]{Example}
\newtheorem{exercise}[subsection]{Exercise}
\newtheorem{situation}[subsection]{Situation}

\theoremstyle{remark}
\newtheorem{remark}[subsection]{Remark}
\newtheorem{remarks}[subsection]{Remarks}

\numberwithin{equation}{subsection}

% Macros
%
\def\lim{\mathop{\mathrm{lim}}\nolimits}
\def\colim{\mathop{\mathrm{colim}}\nolimits}
\def\Spec{\mathop{\mathrm{Spec}}}
\def\Hom{\mathop{\mathrm{Hom}}\nolimits}
\def\Ext{\mathop{\mathrm{Ext}}\nolimits}
\def\SheafHom{\mathop{\mathcal{H}\!\mathit{om}}\nolimits}
\def\SheafExt{\mathop{\mathcal{E}\!\mathit{xt}}\nolimits}
\def\Sch{\mathit{Sch}}
\def\Mor{\mathop{\mathrm{Mor}}\nolimits}
\def\Ob{\mathop{\mathrm{Ob}}\nolimits}
\def\Sh{\mathop{\mathit{Sh}}\nolimits}
\def\NL{\mathop{N\!L}\nolimits}
\def\CH{\mathop{\mathrm{CH}}\nolimits}
\def\proetale{{pro\text{-}\acute{e}tale}}
\def\etale{{\acute{e}tale}}
\def\QCoh{\mathit{QCoh}}
\def\Ker{\mathop{\mathrm{Ker}}}
\def\Im{\mathop{\mathrm{Im}}}
\def\Coker{\mathop{\mathrm{Coker}}}
\def\Coim{\mathop{\mathrm{Coim}}}

% Boxtimes
%
\DeclareMathSymbol{\boxtimes}{\mathbin}{AMSa}{"02}

%
% Macros for moduli stacks/spaces
%
\def\QCohstack{\mathcal{QC}\!\mathit{oh}}
\def\Cohstack{\mathcal{C}\!\mathit{oh}}
\def\Spacesstack{\mathcal{S}\!\mathit{paces}}
\def\Quotfunctor{\mathrm{Quot}}
\def\Hilbfunctor{\mathrm{Hilb}}
\def\Curvesstack{\mathcal{C}\!\mathit{urves}}
\def\Polarizedstack{\mathcal{P}\!\mathit{olarized}}
\def\Complexesstack{\mathcal{C}\!\mathit{omplexes}}
% \Pic is the operator that assigns to X its picard group, usage \Pic(X)
% \Picardstack_{X/B} denotes the Picard stack of X over B
% \Picardfunctor_{X/B} denotes the Picard functor of X over B
\def\Pic{\mathop{\mathrm{Pic}}\nolimits}
\def\Picardstack{\mathcal{P}\!\mathit{ic}}
\def\Picardfunctor{\mathrm{Pic}}
\def\Deformationcategory{\mathcal{D}\!\mathit{ef}}
\usepackage{cleveref}
\crefname{originallemma}{original lemma}{original lemmas}
\Crefname{originallemma}{Original lemma}{Original lemmas}
\makeatletter
\newlabel{criteria-lemma-BG-algebraic}{{06PL}{source}{}{sourcepointer-1}{}}
\newlabel{criteria-lemma-BG-algebraic@cref}{{[originallemma][0][]06PL}{[1][0][]source}{}{}{}}
\makeatother
\crefname{originalremark}{original remark}{original remarks}
\Crefname{originalremark}{Original remark}{Original remarks}
\makeatletter
\newlabel{examples-stacks-remark-X-mod-G-group}{{0371}{source}{}{sourcepointer-2}{}}
\newlabel{examples-stacks-remark-X-mod-G-group@cref}{{[originalremark][0][]0371}{[1][0][]source}{}{}{}}
\makeatother
\crefname{originallemma}{original lemma}{original lemmas}
\Crefname{originallemma}{Original lemma}{Original lemmas}
\makeatletter
\newlabel{perfect-lemma-affine-compare-bounded}{{06Z0}{source}{}{sourcepointer-3}{}}
\newlabel{perfect-lemma-affine-compare-bounded@cref}{{[originallemma][0][]06Z0}{[1][0][]source}{}{}{}}
\makeatother
\crefname{originaldefinition}{original definition}{original definitions}
\Crefname{originaldefinition}{Original definition}{Original definitions}
\makeatletter
\newlabel{stacks-cohomology-definition-flat-base-change}{{0762}{source}{}{sourcepointer-4}{}}
\newlabel{stacks-cohomology-definition-flat-base-change@cref}{{[originaldefinition][0][]0762}{[1][0][]source}{}{}{}}
\makeatother
\crefname{originalproposition}{original proposition}{original propositions}
\Crefname{originalproposition}{Original proposition}{Original propositions}
\makeatletter
\newlabel{stacks-cohomology-proposition-smooth-covering-compute-cohomology}{{0784}{source}{}{sourcepointer-5}{}}
\newlabel{stacks-cohomology-proposition-smooth-covering-compute-cohomology@cref}{{[originalproposition][0][]0784}{[1][0][]source}{}{}{}}
\makeatother
\crefname{originalproposition}{original proposition}{original propositions}
\Crefname{originalproposition}{Original proposition}{Original propositions}
\makeatletter
\newlabel{stacks-perfect-proposition-derived-direct-image-quasi-coherent}{{07BC}{source}{}{sourcepointer-6}{}}
\newlabel{stacks-perfect-proposition-derived-direct-image-quasi-coherent@cref}{{[originalproposition][0][]07BC}{[1][0][]source}{}{}{}}
\makeatother
\crefname{originaldefinition}{original definition}{original definitions}
\Crefname{originaldefinition}{Original definition}{Original definitions}
\makeatletter
\newlabel{stacks-sheaves-definition-locally-quasi-coherent}{{06WJ}{source}{}{sourcepointer-7}{}}
\newlabel{stacks-sheaves-definition-locally-quasi-coherent@cref}{{[originaldefinition][0][]06WJ}{[1][0][]source}{}{}{}}
\makeatother
\crefname{originallemma}{original lemma}{original lemmas}
\Crefname{originallemma}{Original lemma}{Original lemmas}
\makeatletter
\newlabel{stacks-sheaves-lemma-quasi-coherent}{{06WK}{source}{}{sourcepointer-8}{}}
\newlabel{stacks-sheaves-lemma-quasi-coherent@cref}{{[originallemma][0][]06WK}{[1][0][]source}{}{}{}}
\makeatother
\title{A quasi-compact quasi-separated morphism of algebraic stacks that does not preserve D\_QCoh under derived pushforward}
\author{The Stacks Project Authors}
\date{stacks-project:examples.tex:local-snapshot-20261012; GNU FDL 1.2 or later}
\begin{document}
\maketitle
\noindent\textit{Editorial note.} All delivered mathematical text is the author's own, copied from the pinned Stacks Project chapter file examples.tex; no statement or proof step was rewritten, reordered or removed.\par\medskip
\noindent\textit{Editorial note.} The counted proof is the author's complete argument: The middle context span is the statement environment of lemma-is-limit only; its proof is not part of the delivery. The two 'insert future reference here' placeholders flag standard group-cohomology references and sit in the construction layer.\par\medskip
\noindent\textit{Editorial note.} The label lemma-push-not-OK names the counted statement, and the proof environment immediately following it in the source is delivered verbatim as its proof.\par\medskip
\noindent\textit{Editorial note.} The counted statement is carried by the author's own lemma environment in the chapter source.\par\medskip
\noindent\textit{Allowed standard prerequisites (editorial).}\begin{itemize}
\item Context span 4920-4986: k of characteristic p > 0, S = Spec(k[x]), X = [S/G] with G = Z/pZ acting trivially (Examples of Stacks, X-mod-G-group; Criteria for Representability, BG-algebraic); f: X -> S is quasi-compact and quasi-separated so Rf\_\{QCoh,*\}: D\textasciicircum{}+\_QCoh(X) -> D\textasciicircum{}+\_QCoh(S) exists (Derived Categories of Stacks, derived-direct-image-quasi-coherent); RGamma(X, O\_X) is computed by the Cech complex k[x] -> Map(G, k[x]) -> Map(G\textasciicircum{}2, k[x]) -> ... (Cohomology of Stacks, smooth-covering-compute-cohomology; S x\_X ... x\_X S = S x G\textasciicircum{}p affine); this is the group cohomology of G acting trivially on k[x] (the author's two 'insert future reference here' placeholders flag the standard references, kept verbatim per the L48293 precedent); the cohomology of the cyclic group G on k[x] is one copy of k[x] in each degree >= 0, giving Rf\_*O\_X = direct sum over n >= 0 of O\_S[-n]; E = direct sum over m >= 0 of O\_X[m] is an object of D\_QCoh(O\_X)
\item Context span 4989-5005 (lemma-is-limit statement): for K with locally quasi-coherent cohomology sheaves satisfying the flat base change property there is a distinguished triangle K -> prod tau\_\{>= -n\} K -> prod tau\_\{>= -n\} K -> K[1], i.e., K is the derived limit of its canonical truncations; the delivered card L48295 (sum-is-product) gives direct sum = product in D\_QCoh for such K
\item Context span 5041-5057: since a quasi-coherent module is locally quasi-coherent with the flat base change property (Sheaves on Stacks, quasi-coherent) we get E = prod over m >= 0 of O\_X[m]; since cohomology commutes with limits, Rf\_*E = prod over m >= 0 of (direct sum over n >= 0 of O\_S[m-n]); the degree-0 cohomology sheaf is an infinite product of copies of O\_S, which is not even locally quasi-coherent, so this complex is not an object of D\_QCoh(O\_S)
\end{itemize}\medskip
\ifdefined\hypertarget\hypertarget{sourcepointer-1}{}\fi
\noindent\textit{Source reference.} Original lemma 06PL: sources/stacks-chow/tags.txt, source lines 8575--8575. Prerequisite pointer invoked at 5077; it names a standard fact of a sibling chapter that the counted argument uses, and every step of the counted proof built on it is retained verbatim in the delivered body.\par\medskip
\ifdefined\hypertarget\hypertarget{sourcepointer-2}{}\fi
\noindent\textit{Source reference.} Original remark 0371: sources/stacks-chow/tags.txt, source lines 4184--4184. Prerequisite pointer invoked at 5075; it names a standard fact of a sibling chapter that the counted argument uses, and every step of the counted proof built on it is retained verbatim in the delivered body.\par\medskip
\ifdefined\hypertarget\hypertarget{sourcepointer-3}{}\fi
\noindent\textit{Source reference.} Original lemma 06Z0: sources/stacks-chow/tags.txt, source lines 8904--8904. Prerequisite pointer invoked at 5096; it names a standard fact of a sibling chapter that the counted argument uses, and every step of the counted proof built on it is retained verbatim in the delivered body.\par\medskip
\ifdefined\hypertarget\hypertarget{sourcepointer-4}{}\fi
\noindent\textit{Source reference.} Original definition 0762: sources/stacks-chow/tags.txt, source lines 9151--9151. Prerequisite pointer invoked at 5145; it names a standard fact of a sibling chapter that the counted argument uses, and every step of the counted proof built on it is retained verbatim in the delivered body.\par\medskip
\ifdefined\hypertarget\hypertarget{sourcepointer-5}{}\fi
\noindent\textit{Source reference.} Original proposition 0784: sources/stacks-chow/tags.txt, source lines 9224--9224. Prerequisite pointer invoked at 5107; it names a standard fact of a sibling chapter that the counted argument uses, and every step of the counted proof built on it is retained verbatim in the delivered body.\par\medskip
\ifdefined\hypertarget\hypertarget{sourcepointer-6}{}\fi
\noindent\textit{Source reference.} Original proposition 07BC: sources/stacks-chow/tags.txt, source lines 9337--9337. Prerequisite pointer invoked at 5091; it names a standard fact of a sibling chapter that the counted argument uses, and every step of the counted proof built on it is retained verbatim in the delivered body.\par\medskip
\ifdefined\hypertarget\hypertarget{sourcepointer-7}{}\fi
\noindent\textit{Source reference.} Original definition 06WJ: sources/stacks-chow/tags.txt, source lines 8818--8818. Prerequisite pointer invoked at 5143; it names a standard fact of a sibling chapter that the counted argument uses, and every step of the counted proof built on it is retained verbatim in the delivered body.\par\medskip
\ifdefined\hypertarget\hypertarget{sourcepointer-8}{}\fi
\noindent\textit{Source reference.} Original lemma 06WK: sources/stacks-chow/tags.txt, source lines 8819--8819. Prerequisite pointer invoked at 5158; it names a standard fact of a sibling chapter that the counted argument uses, and every step of the counted proof built on it is retained verbatim in the delivered body.\par\medskip
Let $k$ be a field of characteristic $p > 0$. Let $S = \Spec(k[x])$.
Let $G = \mathbf{Z}/p\mathbf{Z}$ viewed either as an abstract group
or as a constant group scheme over $S$. Consider the algebraic stack
$\mathcal{X} = [S/G]$ where $G$ acts trivially on $S$, see
Examples of Stacks, Remark \ref{examples-stacks-remark-X-mod-G-group}
and
Criteria for Representability, Lemma \ref{criteria-lemma-BG-algebraic}.
Consider the structure morphism
$$
f : \mathcal{X} \longrightarrow S
$$
This morphism is quasi-compact and quasi-separated. Hence we get a functor
$$
Rf_{\QCoh, *} :
D^{+}_\QCoh(\mathcal{O}_\mathcal{X})
\longrightarrow
D^{+}_\QCoh(\mathcal{O}_S),
$$
see 
Derived Categories of Stacks, Proposition
\ref{stacks-perfect-proposition-derived-direct-image-quasi-coherent}.
Let's compute $Rf_{\QCoh, *}\mathcal{O}_\mathcal{X}$.
Since $D_\QCoh(\mathcal{O}_S)$ is equivalent to
the derived category of $k[x]$-modules (see
Derived Categories of Schemes, Lemma
\ref{perfect-lemma-affine-compare-bounded})
this is equivalent to computing
$R\Gamma(\mathcal{X}, \mathcal{O}_\mathcal{X})$.
For this we can use the covering $S \to \mathcal{X}$ and the spectral
sequence
$$
H^q(S \times_\mathcal{X} \ldots \times_\mathcal{X} S, O)
\Rightarrow H^{p + q}(\mathcal{X}, \mathcal{O}_\mathcal{X})
$$
see
Cohomology of Stacks, Proposition
\ref{stacks-cohomology-proposition-smooth-covering-compute-cohomology}.
Note that
$$
S \times_\mathcal{X} \ldots \times_\mathcal{X} S = S \times G^p
$$
which is affine. Thus the complex
$$
k[x] \to \text{Map}(G, k[x]) \to \text{Map}(G^2, k[x]) \to \ldots
$$
computes $R\Gamma(\mathcal{X}, \mathcal{O}_\mathcal{X})$.
Here for $\varphi \in \text{Map}(G^{p - 1}, k[x])$ its differential is
the map which sends $(g_1, \ldots, g_p)$ to
$$
\varphi(g_2, \ldots, g_p) +
\sum\nolimits_{i = 1}^{p - 1}
(-1)^i\varphi(g_1, \ldots, g_i + g_{i + 1}, \ldots, g_p)
+ (-1)^p\varphi(g_1, \ldots, g_{p - 1}).
$$
This is just the complex computing the group cohomology of $G$ acting
trivially on $k[x]$ (insert future reference here). The cohomology of
the cyclic group $G$ on $k[x]$ is exactly one copy of $k[x]$ in each
cohomological degree $\geq 0$ (insert future reference here). We conclude
that
$$
Rf_*\mathcal{O}_\mathcal{X} = \bigoplus\nolimits_{n \geq 0} \mathcal{O}_S[-n]
$$
Now, consider the complex
$$
E = \bigoplus\nolimits_{m \geq 0} \mathcal{O}_\mathcal{X}[m]
$$
This is an object of $D_\QCoh(\mathcal{O}_\mathcal{X})$. We interrupt the

\begin{lemma}
\label{lemma-is-limit}
Let $\mathcal{X}$ be an algebraic stack. Let $K$ be an object of
$D(\mathcal{O}_\mathcal{X})$ whose cohomology sheaves are locally
quasi-coherent (Sheaves on Stacks, Definition
\ref{stacks-sheaves-definition-locally-quasi-coherent})
and satisfy the flat base change property (Cohomology of Stacks,
Definition \ref{stacks-cohomology-definition-flat-base-change}).
Then there exists a distinguished triangle
$$
K \to
\prod\nolimits_{n \geq 0} \tau_{\geq -n} K \to
\prod\nolimits_{n \geq 0} \tau_{\geq -n} K \to K[1]
$$
in $D(\mathcal{O}_\mathcal{X})$. In other words, $K$ is the derived
limit of its canonical truncations.
\end{lemma}

\noindent
We continue our discussion. Since a quasi-coherent module is locally
quasi-coherent and satisfies the flat base change property
(Sheaves on Stacks, Lemma \ref{stacks-sheaves-lemma-quasi-coherent}) we get
$$
E = \prod\nolimits_{m \geq 0} \mathcal{O}_\mathcal{X}[m]
$$
Since cohomology commutes
with limits we see that
$$
Rf_*E = \prod\nolimits_{m \geq 0}
\left(\bigoplus\nolimits_{n \geq 0} \mathcal{O}_S[m - n]\right)
$$
Note that this complex is not an object of $D_\QCoh(\mathcal{O}_S)$
because the cohomology sheaf in degree $0$ is an infinite product of copies
of $\mathcal{O}_S$ which is not even a locally quasi-coherent
$\mathcal{O}_S$-module.

\begin{lemma}
\label{lemma-push-not-OK}
A quasi-compact and quasi-separated morphism
$f : \mathcal{X} \to \mathcal{Y}$ of algebraic stacks
need not induce a functor
$Rf_* : D_\QCoh(\mathcal{O}_\mathcal{X}) \to
D_\QCoh(\mathcal{O}_\mathcal{Y})$.
\end{lemma}

\begin{proof}
See discussion above.
\end{proof}

\end{document}
